\section{Dirichlet Character}

\subsection{Definition and basic properties}
\begin{definition}[Dirichlet character]\label{def:dirchar}
    Let $m \ge 1$. We say that $\chi : \mathbb{Z} \to \mathbb{C}$ is a Dirichlet character $\pmod m$ if
    \begin{enumerate}[label={\upshape\roman*)}]
        \item\label{def:dirchar:multiplicity}$\chi(ab) = \chi(a)\chi(b), \text{ } \forall a,b \in \mathbb{Z}$;
        \item\label{def:dirchar:zero_condition}$\chi(a) =0 \Leftrightarrow (a,m) \neq1$;
        \item\label{def:dirchar:periodicity}$\chi(a+m) = \chi(a)$, $\forall a \in \mathbb{Z}$.
    \end{enumerate}
\end{definition}
Using \cref{def:dirchar} \ref{def:dirchar:multiplicity}, we obtain the following: $\chi(1)= \chi(1\cdot1) = \chi(1)\chi(1)$ and $\chi(1)\neq0$, which implies $\chi(1)=1$. By periodicity, $\chi(a)$ only depends on the class of $a \pmod m$, so we can see $\chi$ as a function on $\mathbb{Z}/m\mathbb{Z}$.

\begin{definition}[Principal character]
    We define
    $$\chi_0(n):=
    \begin{cases}
        1 & \text{if } (n,m)=1,\\
        0 & \text{if } (n,m) > 1,
    \end{cases}$$
    the principal character $\pmod m$.
\end{definition}
We usually denote a non-principal character with $\psi$.

It is convenient to work in the more general setting of finite abelian groups. Let $G$ be a finite abelian group. A \emph{character} of $G$ is a group homomorphism $\lambda : G\to\mathbb{C}^\times$. We denote by $\widehat G$ the set of characters of $G$.

\begin{lemma}\label{lem:dirchar_group_char}
    The Dirichlet characters $\pmod m$ are in bijection with the characters of $G = \left(\mathbb{Z}/m\mathbb{Z}\right)^\times$: to $\chi$, we associate its restriction $\lambda(a \bmod m) = \chi(a)$ for $(a,m) = 1$; conversely, to $\lambda\in\widehat G$, we associate $\chi(n) = \lambda(n\bmod m)$ if $(n,m) = 1$ and $\chi(n) = 0$ otherwise.
\end{lemma}
\begin{proof}
    If $\chi$ is a Dirichlet character, its restriction is well defined by periodicity, takes non-zero values by \ref{def:dirchar:zero_condition} and is multiplicative by \ref{def:dirchar:multiplicity}, so it is a homomorphism $G\to\mathbb{C}^\times$. Conversely, let $\lambda \in \widehat G$ and $\chi$ defined as above. Conditions \ref{def:dirchar:zero_condition} and \ref{def:dirchar:periodicity} are clear. For \ref{def:dirchar:multiplicity}: if $(a,m) = (b,m) = 1$, then $(ab,m) = 1$ and $\chi(ab) = \lambda(a)\lambda(b) = \chi(a)\chi(b)$; otherwise $(ab,m)\neq1$ and both sides are $0$. The two maps are inverse of each other since a Dirichlet character is determined by its values on the integers coprime to $m$.
\end{proof}

The key fact on characters is that there are \emph{many} of them.
\begin{lemma}[Extension of characters]\label{lem:extension_char}
    Let $G$ be a finite abelian group and $H$ a subgroup of $G$. Every character of $H$ extends to exactly $[G:H]$ characters of $G$.
\end{lemma}
\begin{proof}
    We proceed by induction on $[G:H]$. If $[G:H] = 1$, there is nothing to prove. Otherwise, let $g\in G\setminus H$, let $H' = \langle H,g\rangle$ and let $n\ge 2$ be the smallest integer such that $g^n\in H$ (it exists since $g^{\lvert G\rvert} = 1 \in H$).

    Every element of $H'$ can be written as $hg^j$ with $h\in H$ and $0\le j<n$, and this writing is unique: if $hg^j = h'g^{j'}$ with $0 \le j' \le j < n$, then $g^{j-j'}\in H$, so $j = j'$ by minimality of $n$, and then $h = h'$. In particular $[H':H] = n$.

    Let $\lambda\in\widehat H$. If $\lambda'$ is an extension of $\lambda$ to $H'$, then $w = \lambda'(g)$ satisfies $w^n = \lambda(g^n)$, and $\lambda'(hg^j) = \lambda(h)w^j$, so $\lambda'$ is determined by $w$. Conversely, let $w$ be one of the $n$ complex roots of $w^n = \lambda(g^n)$ (note $\lambda(g^n)\neq0$) and define $\lambda'(hg^j) = \lambda(h)w^j$, which is well defined by uniqueness of the writing. It is a homomorphism: if $j + j' <n$, it is clear, and if $j + j'\ge n$, then $hg^j\cdot h'g^{j'} = (hh'g^n)g^{j+j'-n}$ and
    \begin{equation*}
        \lambda'\left(hh'g^n g^{j+j'-n}\right) = \lambda(hh')\lambda(g^n)w^{j+j'-n} = \lambda(h)\lambda(h')w^{j+j'} = \lambda'(hg^j)\lambda'(h'g^{j'}).
    \end{equation*}
    So $\lambda$ has exactly $n = [H':H]$ extensions to $H'$. Since $[G:H'] < [G:H]$, by induction each of them has exactly $[G:H']$ extensions to $G$, and two different extensions to $H'$ give different extensions to $G$. Every extension of $\lambda$ to $G$ is obtained this way (by restricting it to $H'$). Thus $\lambda$ has $[H':H][G:H'] = [G:H]$ extensions to $G$.
\end{proof}

\begin{proposition}[Basic properties of Dirichlet character]
    \label{prop:number_dir_char}
    There exist $\phi(m)$ distinct Dirichlet characters $\pmod m$ and the set of all distinct characters $\pmod m$ forms a finite abelian group under multiplication, denoted $\widehat{\left(\mathbb{Z}/m \mathbb{Z}\right)^{\times}}$, whose identity is $\chi_0$. The values $\chi(a)$, for $(a,m) = 1$, are $\phi(m)$-th roots of unity, so $\chi^{-1} = \overline{\chi}$. We also have $\widehat{\left(\mathbb{Z}/m \mathbb{Z}\right)^{\times}} \cong\left(\mathbb{Z}/m \mathbb{Z}\right)^{\times}$.
\end{proposition}
\begin{proof}
    Let $G = \left(\mathbb{Z}/m\mathbb{Z}\right)^\times$, so that $\lvert G\rvert = \phi(m)$. By \cref{lem:dirchar_group_char}, it is enough to study $\widehat G$.

    \emph{Number.} The trivial group $H = \{1\}$ has exactly one character, so by \cref{lem:extension_char}, $\lvert\widehat G\rvert = [G:\{1\}] = \phi(m)$.

    \emph{Group structure.} If $\chi,\chi'\in\widehat G$, then $\chi\chi'(a) := \chi(a)\chi'(a)$ is again a homomorphism, the product is associative and commutative, and $\chi_0$ is the identity. For $(a,m)=1$, Euler's theorem (\cref{thm:euler_theorem}) gives $a^{\phi(m)}\equiv1\pmod m$, so $\chi(a)^{\phi(m)} = \chi(a^{\phi(m)}) = \chi(1) = 1$: $\chi(a)$ is a root of unity, hence $\lvert\chi(a)\rvert = 1$ and $\chi(a)^{-1} = \overline{\chi(a)}$. Thus $\overline\chi$ is a character and $\chi\overline\chi = \chi_0$.

    \emph{Isomorphism.} If $C_n = \langle g\rangle$ is cyclic of order $n$, a character is determined by $\lambda(g)$, which can be any $n$-th root of unity; hence $\lambda \mapsto \lambda(g)$ is an isomorphism $\widehat{C_n}\cong\mu_n\cong C_n$. If $G = G_1\times G_2$, then $(\lambda_1,\lambda_2)\mapsto \left((g_1,g_2)\mapsto\lambda_1(g_1)\lambda_2(g_2)\right)$ is an isomorphism $\widehat{G_1}\times\widehat{G_2}\cong \widehat{G}$ (the inverse is given by restriction to $G_1\times\{1\}$ and $\{1\}\times G_2$). By the structure theorem of finite abelian groups, $G$ is a product of cyclic groups, so $\widehat G\cong G$.
\end{proof}
\noindent\emph{Remark.} The isomorphism $\widehat G\cong G$ is not canonical and will not be used. Only the facts that $\lvert\widehat G\rvert = \lvert G\rvert$ and \cref{lem:existence_dir_char} below will matter.

\subsection{Orthogonality relation}
To prove the orthogonality relation we need the following lemma:
\begin{lemma}\label{lem:existence_dir_char}\leavevmode
    \begin{enumerate}[label={\upshape\roman*)}]
        \item\label{lem:existence_dir_char:i} Let $m\ge 2$ and $a \in \mathbb{Z}$ such that $(a,m)=1$ with $a\not \equiv 1 \pmod m$. There exists $\chi \pmod m$ such that $\chi(a) \neq 1$.
        \item\label{lem:existence_dir_char:ii} For $m\ge 2$, let $\chi$ be a character $\pmod m$ with $\chi \neq \chi_0$. Then there exists $a \in \mathbb{Z}$ with $(a,m) = 1$ such that $\chi(a)\neq 1$ (in particular $\chi(a) \neq 0,1$).
    \end{enumerate}
\end{lemma}
\begin{proof}\leavevmode
    \begin{enumerate}[label={\upshape\roman*)}]
        \item Let $G = \left(\mathbb{Z}/m\mathbb{Z}\right)^\times$ and let $H = \langle a\rangle$ be the subgroup generated by the class of $a$. Since $a\not\equiv1\pmod m$, the order $d$ of $a$ in $G$ satisfies $d\ge2$. The map $\lambda(a^j) = e^{2\pi ij/d}$ is a well defined character of $H$ (because $a^j = a^{j'}$ if and only if $d\mid j-j'$), and $\lambda(a) = e^{2\pi i/d}\neq1$. By \cref{lem:extension_char}, $\lambda$ extends to a character of $G$, which gives, by \cref{lem:dirchar_group_char}, a Dirichlet character $\chi$ with $\chi(a) \neq 1$.
        \item The characters $\chi$ and $\chi_0$ are both $0$ on the integers not coprime to $m$. Since $\chi \neq\chi_0$, they must differ at some $a$ with $(a,m) = 1$, that is $\chi(a)\neq\chi_0(a) = 1$. Also $\chi(a) \neq 0$ by \cref{def:dirchar} \ref{def:dirchar:zero_condition}.
    \end{enumerate}
\end{proof}

\begin{theorem}[Orthogonality relation of Dirichlet character]\leavevmode
    \label{thm:orthogonality_rel}
    \begin{enumerate}[label={\upshape\roman*)}]
        \item $\displaystyle\sum_{a \in \mathbb{Z} / m \mathbb{Z}} \chi(a) = \begin{cases}
                \phi(m)  &\text{ if } \chi= \chi_0,\\
                0 &\text{ if } \chi \neq \chi_0;
            \end{cases}$
        \item\label{item:orthogonality_rel_item:ii} $\displaystyle\sum_{\chi \bmod m} \chi(a) = \begin{cases}
                \phi(m) &\text{ if } a \equiv 1 \pmod m, \\
                0 &\text{ if } a \not\equiv 1 \pmod m.
            \end{cases}$
    \end{enumerate}
\end{theorem}

\begin{proof}\leavevmode
    \begin{enumerate}[label={\roman*)}]
        \item\label{1} Since $\chi(a) = 0$ when $(a,m) \neq 1$, we can restrict the sum to $a \in \left(\mathbb{Z} / m \mathbb{Z}\right)^{\times}$. If $\chi = \chi_0$, every term is $1$ and there are $\phi(m)$ of them:
            \begin{equation*}
                \sum_{a \in \mathbb{Z} / m \mathbb{Z}} \chi_0(a)=  \sum_{a \in \mathbb{Z} / m \mathbb{Z}} \mathds{1}_{(a,m)=1} = \phi(m).
            \end{equation*}
            Now suppose $\chi \neq \chi_0$ and, by \cref{lem:existence_dir_char} \ref{lem:existence_dir_char:ii}, take $\alpha$ with $(\alpha,m) = 1$ and $\chi(\alpha) \neq 1$. We have
            \begin{equation}\label{eq:multiplicity_dir_char}
                \chi(\alpha)\sum_{a \in \left(\mathbb{Z} / m \mathbb{Z}\right)^{\times}} \chi(a) = \sum_{a \in \left(\mathbb{Z} / m \mathbb{Z}\right)^{\times}} \chi(\alpha a).
            \end{equation}
            Since $\alpha$ is invertible $\pmod m$, the map $a \mapsto \alpha a$ is a bijection of $\left(\mathbb{Z} / m \mathbb{Z}\right)^{\times}$ onto itself. This implies:
            \begin{equation}\label{eq:complete_set_residu}
                \sum_{a \in \left(\mathbb{Z} / m \mathbb{Z}\right)^{\times}} \chi(\alpha a)= \sum_{a \in \left(\mathbb{Z} / m \mathbb{Z}\right)^{\times}} \chi(a).
            \end{equation}
            Combining \cref{eq:multiplicity_dir_char} and \cref{eq:complete_set_residu}, we get
            \begin{equation*}
                (\chi(\alpha)-1)\sum_{a \in \left(\mathbb{Z} / m \mathbb{Z}\right)^{\times}} \chi( a) = 0.
            \end{equation*}
            Since we took $\alpha$ such that $\chi(\alpha) \neq 1 $, the sum is $0$.
        \item If $a \equiv 1 \pmod m$, then $\chi(a) = \chi(1) = 1$ for every $\chi$ and, by \cref{prop:number_dir_char},
            \begin{equation*}
                \sum_{\chi \bmod m} \chi(a) = \phi(m).
            \end{equation*}
            If $(a,m)\neq1$, then every term is $0$. Now suppose $(a,m) = 1$ and $a\not \equiv 1 \pmod m$. By \cref{lem:existence_dir_char} \ref{lem:existence_dir_char:i}, there exists $\chi'$ such that $\chi'(a) \neq 1$. We use the same idea as in \ref{1}: since $\chi\mapsto\chi'\chi$ is a bijection of the group of characters onto itself,
            \begin{equation*}
                \chi'(a)\sum_{\chi \bmod m}\chi(a) = \sum_{\chi \bmod m} (\chi' \chi)(a)=\sum_{\chi \bmod m}  \chi(a).
            \end{equation*}
            Thus $(\chi'(a)-1)\sum_{\chi}\chi(a) = 0$ and the sum is $0$.
    \end{enumerate}
\end{proof}
Let $\alpha,\beta\in\mathbb{Z}$ with $(\beta,m) = 1$, and let $\beta^{-1}$ be an inverse of $\beta \pmod m$. Applying \ref{item:orthogonality_rel_item:ii} to $a = \alpha \beta^{-1}$, and using $\chi(\beta^{-1}) = \chi(\beta)^{-1} = \overline\chi(\beta)$, we get an equation that will be more useful later on:
\begin{equation}\label{eq:orthogonality_rel_special_case}
    \sum_{\chi \bmod m} \chi(\alpha ) \overline\chi(\beta) =\begin{cases}
                \phi(m) &\text{ if } \alpha \equiv \beta \pmod m, \\
                0 &\text{ if } \alpha \not\equiv \beta \pmod m.
            \end{cases}
\end{equation}
